% Generated draft: requires mathematical and source review.
\documentclass[11pt]{article}
\usepackage[margin=25mm]{geometry}
\usepackage{fontspec}
\setmainfont{Latin Modern Roman}
\usepackage{amsmath,amssymb,mathtools,unicode-math}
\setmathfont{Latin Modern Math}
\usepackage{enumitem,needspace,xurl,hyperref}
\setlength{\parindent}{0pt}
\setlength{\parskip}{5pt}
\setlength{\emergencystretch}{3em}
\newcommand{\N}{\mathbb{N}}
\newcommand{\Z}{\mathbb{Z}}
\newcommand{\Q}{\mathbb{Q}}
\newcommand{\R}{\mathbb{R}}
\newcommand{\C}{\mathbb{C}}
\newcommand{\F}{\mathbb{F}}
\newcommand{\A}{\mathbb{A}}
\newcommand{\PP}{\mathbb{P}}
\newcommand{\E}{\mathbb{E}}
\newcommand{\eps}{\varepsilon}
\newcommand{\dd}{\,\mathrm{d}}
\begin{document}
\section*{UnsolvedMath ID: 700541}
\textbf{Source problem code:} catalogue entry
\subsection{Definitions and mathematical statement}
Work in the usual group language, with multiplication, inversion, and the identity constant. A quasiidentity is a universally quantified formula
\[
\forall x_1\cdots\forall x_n\;\left(\bigwedge_{i=1}^{m}u_i(x_1,\ldots,x_n)=v_i(x_1,\ldots,x_n)\ \Longrightarrow\ u(x_1,\ldots,x_n)=v(x_1,\ldots,x_n)\right),
\]
where all displayed terms are group terms and the empty conjunction, when \(m=0\), is true. For a group \(G\), let \(\operatorname{QId}(G)\) be the set of all quasiidentities valid in \(G\).

A set \(B\subseteq\operatorname{QId}(G)\) is a basis of this theory if, for every group \(H\),
\[
H\models B\quad\Longleftrightarrow\quad H\models\operatorname{QId}(G).
\]
The basis is independent if, for every \(\varphi\in B\), there exists a group \(H_\varphi\) such that
\[
H_\varphi\models B\setminus\{\varphi\},\qquad H_\varphi\not\models\varphi.
\]
Thus consequence and independence are taken relative to the axioms of groups, with arbitrary groups as models.

\textbf{Question (V. A. Gorbunov, 1984; Kourovka Notebook 9.6).} Is it true that for every finite group \(G\), every independent basis \(B\) of \(\operatorname{QId}(G)\) is finite?

This formulation reads ``an independent basis'' as quantifying over every independent basis that exists. It does not assert that every finite group has an independent basis. The supplied record does not explicitly distinguish this reading from a question about existence of a finite independent basis, and does not specify its conventions for logical consequence; these interpretations require source review.
\subsection{Short English statement}
For each finite group, must every independent collection of quasiidentities that axiomatizes all its quasiidentities be finite? Independence means that removing any member loses a consequence, with consequence evaluated among all groups. This reading does not assume that an independent basis exists.
\subsection{Sources}
\begin{itemize}
\item Local UnsolvedMath registry, supplied complete source record: permanent ID \(700541\), exact source problem code \texttt{catalogue entry}, title \emph{Kourovka Notebook Problem 9.6}. The registry labels the problem \texttt{open}; this is supplied classification, not an independently verified mathematical or literature assessment.
\item Supplied bibliographic attribution: \emph{The Kourovka Notebook (Unsolved Problems in Group Theory)}, original problem 9.6, ninth issue (1984); proposed by V. A. Gorbunov. The record also identifies the twenty-first issue edition (2026), arXiv:1401.0300v48. Its collection description names E. I. Khukhro and V. D. Mazurov as editors.
\item The registry record was created and updated on \texttt{2026-10-16T14:00:00Z}, with upstream revision \texttt{d371ee917ac9647fbb06cd6fc529deec6f3dda82}. These are record provenance dates, not dates of independent source verification. The cited publication has not been independently consulted for this draft.
\end{itemize}
\end{document}
